\documentclass[12pt,a4paper]{article}

\usepackage[a4paper,margin=2cm]{geometry}
\usepackage[T2A]{fontenc}
\usepackage[utf8]{inputenc}
\usepackage[russian]{babel}
\usepackage{amsmath,amssymb,mathtools}
\usepackage{array,booktabs,longtable}
\usepackage{xcolor}
\usepackage{hyperref}
\usepackage{tikz}
\usepackage{fancyhdr}

\hypersetup{
  colorlinks=true,
  linkcolor=blue!55!black,
  urlcolor=blue!55!black
}

\setlength{\parindent}{0pt}
\setlength{\parskip}{0.55em}
\renewcommand{\arraystretch}{1.25}

\pagestyle{fancy}
\fancyhf{}
\fancyfoot[C]{\small Лёвин Артём Александрович эксклюзивно для Ксении Васильченко}
\renewcommand{\headrulewidth}{0pt}
\renewcommand{\footrulewidth}{0pt}

\newcommand{\err}{\textcolor{red}{Ошибка:}}
\newcommand{\idea}{\textcolor{blue!60!black}{Ключевая идея:}}
\newcommand{\result}{\textcolor{teal!60!black}{Итог:}}

\begin{document}

\thispagestyle{empty}
\vspace*{\fill}
\begin{center}
{\LARGE\bfseries Ревью домашней работы\par}
\vspace{1.2cm}
{\large Автор: Лёвин Артём Александрович\par}
\vspace{0.7cm}
{\large 3 октября 2026 года\par}
\end{center}
\vspace*{\fill}
\begin{center}
\begin{tikzpicture}
\draw[blue!35!black,line width=0.7pt]
(0,0) .. controls (2,0.7) and (4,-0.7) .. (6,0);
\end{tikzpicture}
\end{center}

\newpage

{\small
\setlength{\tabcolsep}{3pt}
\begin{longtable}{>{\raggedright\arraybackslash}p{0.10\linewidth}
                  >{\raggedright\arraybackslash}p{0.19\linewidth}
                  >{\raggedright\arraybackslash}p{0.29\linewidth}
                  >{\raggedright\arraybackslash}p{0.17\linewidth}
                  >{\raggedright\arraybackslash}p{0.17\linewidth}}
\toprule
№ задания & Тема & Суть ошибки & Ваш ответ & Правильный ответ \\
\midrule
\endfirsthead
\toprule
№ задания & Тема & Суть ошибки & Ваш ответ & Правильный ответ \\
\midrule
\endhead

\hyperref[task:1]{1} &
Признаки делимости &
В пункте о делимости на четыре включено число 726, хотя его последние две цифры образуют число 26, которое на четыре не делится. &
В пункте в): 144, 726, 1248 &
В пункте в): 144, 1248 \\
\midrule

\hyperref[task:5]{5} &
Взаимно простые числа &
Ответы «да» и «нет» получены верно, но по условию требовалось кратко обосновать каждый ответ. &
а) да; б) нет &
а) взаимно простые; б) не взаимно простые, общий делитель 7 \\
\midrule

\hyperref[task:6]{6} &
Чётность выражений &
Чётность обоих выражений определена верно, но отсутствуют обязательные объяснения без вычисления самих значений. &
а) нечётное; б) нечётное &
Оба выражения нечётные; требуется обоснование по правилам чётности \\
\bottomrule
\end{longtable}
}

\newpage
\thispagestyle{empty}
\vspace*{\fill}
\begin{center}
{\LARGE\bfseries Разбор ошибок}
\end{center}
\vspace*{\fill}

\newpage
\phantomsection\label{task:1}
\section*{Задание №1}

\textbf{Текст задания}

Даны числа
\[
144,\ 315,\ 726,\ 935,\ 1248,\ 1725.
\]
Нужно выписать среди них все числа, которые делятся на два, на три, на четыре, на пять и на шесть.

\textbf{Ваше решение}

В пунктах о делимости на два, на три, на пять и на шесть списки составлены правильно. В пункте о делимости на четыре записаны числа
\[
144,\ 726,\ 1248.
\]

\textbf{Ваш ответ}

В пункте в):
\[
144,\ 726,\ 1248.
\]

\textbf{\err\ в список чисел, делящихся на четыре, ошибочно включено число 726.}

\textbf{Где именно возникает ошибка}

Последний правильный шаг --- включение числа 144 в список: его последние две цифры образуют число 44, а 44 делится на четыре.

Первый неверный шаг --- включение числа 726 в этот же список.

\textbf{Почему этот шаг неверен}

Признак делимости на четыре формулируется так: натуральное число делится на четыре тогда и только тогда, когда число, образованное его последними двумя цифрами, делится на четыре.

У числа 726 последние две цифры --- 26. Число 26 на четыре не делится, потому что ближайшие кратные четырём --- 24 и 28.

Следовательно, число 726 не делится на четыре и не должно находиться в ответе этого пункта.

Для числа 1248 последние две цифры --- 48. Число 48 делится на четыре, поэтому 1248 включено правильно.

\textbf{\idea} при проверке делимости на четыре не нужно проверять всё число целиком: достаточно посмотреть только на последние две цифры.

\textbf{Исправление от последнего правильного шага}

Проверяем числа по последним двум цифрам:
\[
144:\ 44 \text{ делится на }4;
\]
\[
315:\ 15 \text{ не делится на }4;
\]
\[
726:\ 26 \text{ не делится на }4;
\]
\[
935:\ 35 \text{ не делится на }4;
\]
\[
1248:\ 48 \text{ делится на }4;
\]
\[
1725:\ 25 \text{ не делится на }4.
\]

Поэтому в пункте о делимости на четыре должны остаться только два числа.

\textbf{Полное правильное решение}

\textbf{а) Делятся на два.}

Число делится на два, если его последняя цифра чётная. Получаем:
\[
144,\ 726,\ 1248.
\]

\textbf{б) Делятся на три.}

Число делится на три, если сумма его цифр делится на три.

Для 144 сумма цифр равна девяти; для 315 --- девяти; для 726 --- пятнадцати; для 1248 --- пятнадцати; для 1725 --- пятнадцати. Все эти суммы делятся на три.

Получаем:
\[
144,\ 315,\ 726,\ 1248,\ 1725.
\]

\textbf{в) Делятся на четыре.}

Проверяем последние две цифры. На четыре делятся только 44 и 48, поэтому:
\[
144,\ 1248.
\]

\textbf{г) Делятся на пять.}

Число делится на пять, если оканчивается на ноль или на пять. Получаем:
\[
315,\ 935,\ 1725.
\]

\textbf{д) Делятся на шесть.}

Число должно одновременно делиться на два и на три. Получаем:
\[
144,\ 726,\ 1248.
\]

\textbf{Проверка}

Для ошибочно включённого числа 726:
\[
726=4\cdot181+2.
\]
Остаток не равен нулю, следовательно, число действительно не делится на четыре.

\textbf{\result} правильный ответ пункта в):
\[
\boxed{144,\ 1248.}
\]

\textbf{Задача на закрепление}

Среди чисел
\[
132,\ 258,\ 420,\ 735,\ 1024,\ 1416
\]
выпишите все числа, которые делятся на четыре.

\newpage
\phantomsection\label{task:5}
\section*{Задание №5}

\textbf{Текст задания}

Не ограничиваясь ответом «да» или «нет», определить, какие пары чисел являются взаимно простыми:

а) 14 и 25;

б) 21 и 49.

Кратко обосновать каждый ответ.

\textbf{Ваше решение}

Записано:

а) да;

б) нет.

\textbf{Ваш ответ}

а) да;

б) нет.

\textbf{\err\ математический результат верный, но отсутствует обязательное обоснование.}

\textbf{Где именно возникает ошибка}

Последний правильный шаг --- верная классификация обеих пар.

Первый неполный шаг --- завершение решения словами «да» и «нет», хотя условие прямо требует объяснить каждый ответ.

\textbf{Почему это неполно}

Два натуральных числа называются взаимно простыми, если их наибольший общий делитель равен единице.

Поэтому для каждой пары нужно показать, почему общего делителя больше единицы нет или, наоборот, предъявить общий делитель больше единицы.

Для первой пары удобно разложить числа на простые множители:
\[
14=2\cdot7,
\qquad
25=5^2.
\]
Общих простых множителей нет, значит, их наибольший общий делитель равен единице.

Для второй пары:
\[
21=3\cdot7,
\qquad
49=7^2.
\]
У обоих чисел есть общий множитель семь, поэтому эта пара взаимно простой не является.

\textbf{\idea} чтобы обосновать взаимную простоту, достаточно показать, что наибольший общий делитель равен единице; чтобы опровергнуть её, достаточно найти общий делитель больше единицы.

\textbf{Исправление}

Для первой пары записываем: числа 14 и 25 взаимно простые, потому что у них нет общих простых множителей и их наибольший общий делитель равен единице.

Для второй пары записываем: числа 21 и 49 не являются взаимно простыми, потому что оба делятся на семь.

\textbf{Полное правильное решение}

\textbf{а) Числа 14 и 25.}

Разложим их:
\[
14=2\cdot7,
\qquad
25=5^2.
\]
Общих простых множителей нет. Значит,
\[
\text{НОД}(14,25)=1.
\]
Следовательно, числа 14 и 25 взаимно простые.

\textbf{б) Числа 21 и 49.}

Разложим их:
\[
21=3\cdot7,
\qquad
49=7^2.
\]
Оба числа имеют общий множитель семь. Поэтому
\[
\text{НОД}(21,49)=7.
\]
Наибольший общий делитель больше единицы, следовательно, числа 21 и 49 не являются взаимно простыми.

\textbf{Проверка}

Первая пара не имеет общего простого делителя. Во второй паре число семь делит оба числа без остатка. Это подтверждает оба вывода.

\textbf{\result} а) 14 и 25 взаимно простые; б) 21 и 49 не взаимно простые.

\textbf{Задача на закрепление}

Определите, являются ли взаимно простыми пары 16 и 27, а также 35 и 49. Каждый ответ кратко обоснуйте.

\newpage
\phantomsection\label{task:6}
\section*{Задание №6}

\textbf{Текст задания}

Не выполняя сами вычисления, определить, чётным или нечётным будет значение выражения.

а)
\[
15\cdot17+24\cdot31;
\]

б)
\[
13+15+17+19+21.
\]

Для каждого пункта записать короткое объяснение.

\textbf{Ваше решение}

Записано:

а) нечётное;

б) нечётное.

\textbf{Ваш ответ}

а) нечётное;

б) нечётное.

\textbf{\err\ ответы верные, но отсутствует требуемое объяснение по правилам чётности.}

\textbf{Где именно возникает ошибка}

Последний правильный шаг --- верное определение чётности обоих выражений.

Первый неполный шаг --- отсутствие объяснения, почему каждое выражение имеет именно такую чётность.

\textbf{Почему это неполно}

В условии специально сказано не выполнять вычисления. Значит, нужно использовать свойства чётных и нечётных чисел.

Произведение двух нечётных чисел нечётно. Если хотя бы один множитель чётный, произведение чётно. Сумма нечётного и чётного числа нечётна.

Кроме того, сумма нечётного количества нечётных слагаемых нечётна.

Именно эти правила позволяют получить ответ без вычисления самих значений выражений.

\textbf{\idea} определяем только тип каждого множителя и слагаемого --- чётный или нечётный --- и затем применяем правила для суммы и произведения.

\textbf{Исправление}

В пункте а) числа 15 и 17 нечётные, поэтому их произведение нечётное. Число 24 чётное, поэтому произведение 24 и 31 чётное. Сумма нечётного и чётного числа нечётна.

В пункте б) складываются пять нечётных чисел. Пять --- нечётное количество слагаемых, поэтому их сумма нечётна.

\textbf{Полное правильное решение}

\textbf{а)}

Числа 15 и 17 нечётные. Значит, их произведение нечётное.

Число 24 чётное. Значит, произведение 24 и 31 чётное независимо от чётности второго множителя.

Получаем сумму нечётного и чётного числа. Такая сумма нечётна.

Следовательно, значение первого выражения нечётное.

\textbf{б)}

Все пять слагаемых
\[
13,\ 15,\ 17,\ 19,\ 21
\]
нечётные.

Сумма двух нечётных чисел чётна. Сумма следующих двух нечётных чисел тоже чётна. После этого остаётся ещё одно нечётное слагаемое. Сумма двух чётных чисел и одного нечётного числа нечётна.

Следовательно, значение второго выражения также нечётное.

\textbf{Проверка}

В первом пункте структура чётности имеет вид
\[
\text{нечётное}+\text{чётное}=\text{нечётное}.
\]

Во втором пункте имеется нечётное количество нечётных слагаемых, поэтому итоговая сумма нечётна.

\textbf{\result} оба выражения имеют нечётные значения.

\textbf{Задача на закрепление}

Не выполняя вычислений, определите чётность выражения
\[
17\cdot21+32\cdot15
\]
и кратко объясните ответ.

\end{document}
